\clearpage\fancyhead[L]{\small KOMET.}\SourceLeaf{346}
\begin{center}\textbf{KOMET.}\par\vspace{4pt}\hrule\vspace{4pt}\fontsize{10}{12}\selectfont Monatliche Correspondenz. Band XXIV, S. 507\,..\,517, November 1811.\par\vspace{4pt}\hrule\vspace{4pt}\end{center}

\begin{flushright}Göttingen, 15. November.\end{flushright}
\par\indent Nach dem, was \textsc{Pingré} aus europäischen und chinesischen Beobachtungen des Kometen von 1301 anführt, \emph{kann dieser Komet mit dem gegenwärtigen nicht identisch sein.} Denn der gegenwärtige Komet kann unmöglich am 16. September bei einer Länge von 110$^\circ$ erschienen sein. Ebenso unmöglich ist es, dass er am 30. September bei einer Länge von 231$^\circ$ und \emph{nur} 26$^\circ$ nördlicher Breite erschienen ist.


\par\indent Teils wegen ungünstigen Wetters, teils wegen anderer Verpflichtungen konnte ich seit meinem letzten Brief nur eine Beobachtung des Kometen machen. Ich teile sie Ihnen hier mit.

\begin{center}
\fontsize{10}{12}\selectfont\setlength{\tabcolsep}{3pt}\renewcommand{\arraystretch}{1.15}\begin{tabular}{l|c|c}
 & Rektasz. & Nördl. Dekl. \\ \hline
1811 Okt. 20 $10^{\mathrm u}$ $9^{\mathrm m}$ $12^{\mathrm s}$ m. Z. & $247^\circ\,46'\,39''$ & $41^\circ\,20'\,15''$ \\
\end{tabular}
\end{center}


\par\indent Ich habe diese Beobachtungen mit meinen letzten parabolischen Elementen verglichen, ebenso wie Herr \textsc{Nicolai} die letzten Beobachtungen von \textsc{Olbers} und die letzten Seeberger Beobachtungen. Es folgen die Ergebnisse.

\Needspace{20\baselineskip}\begin{center}Bremer Beobachtungen\end{center}
\begin{center}\fontsize{10}{12}\selectfont\setlength{\tabcolsep}{8pt}\renewcommand{\arraystretch}{1.15}\begin{tabular}{l|c|c}
& \multicolumn{2}{c}{Unterschied} \\
Tag der & & \\
Beobachtung & in Rektasz. & in Dekl. \\ \hline
Okt. 19 & $-73''$ 
 & $+152''$ \\
19 & $-73$ & ---\;--- \\
20 & $-39$ & $+139$ \\
24 & $-82$ & $+227$ \\
24 & $-84$ & $+205$ \\
25 & $-108$ & ---\;--- \\
25 & $-91$ & ---\;--- \\
28 & $-60$ & $+293$ \\
Nov. 4 & $-159$ & ---\;--- \\
4 & $-155$ & $+342$ \\
7 & $-91$ & ---\;--- \\
7 & ---\;--- & $+377$ \\
7 & $-102$ & ---\;--- \\
9 & $-274$ & $+379$ \\
\end{tabular}\end{center}
\Needspace{15\baselineskip}\begin{center}Seeberger Beobachtungen\end{center}
\begin{center}\fontsize{10}{12}\selectfont\setlength{\tabcolsep}{8pt}\renewcommand{\arraystretch}{1.15}\begin{tabular}{l|c|c}
& \multicolumn{2}{c}{Unterschied} \\
Tag der & & \\
Beobachtung & in Rektasz. & in Dekl. \\ \hline
Okt. 20 & $-113''$ & $+196''$ \\
24 & $-109$ & $+175$ \\
25 & $-88$ & $+263$ \\
28 & $-87$ & $+327$ \\
29 & $-124$ & $+265$ \\
Nov. 4 & $-110$ & $+387$ \\
5 & $-74$ & ---\;--- \\
\end{tabular}\end{center}


\Needspace{19\baselineskip}\par\indent Dr. \textsc{Olbers} hat mir die ersten Beobachtungen von Herrn \textsc{Flaugergues} mitgeteilt. Ich füge Herrn \textsc{Nicolai}s Berechnung der Abweichungen von den Elementen bei

\Needspace{15\baselineskip}
\begin{center}\fontsize{9}{11}\selectfont\setlength{\tabcolsep}{3pt}\renewcommand{\arraystretch}{1.15}\begin{tabular}{l|c|c|c|c}
\shortstack{Tag der\\Beobachtung} & \shortstack{Rektasz.\\des Kometen} & \shortstack{Abweichung von\\den Elementen} & \shortstack{Südl. Dekl.\\des Kometen} & \shortstack{Abweichung von\\den Elementen}\\\hline
März 26 & $120^\circ\,16'\,0''$ & $+859''$ & $29^\circ\,15'\,0''$ & $+300''$\\
28 & $119\;52\;56$ & $+316$ & $28\;\;7\;\;0$ & $+257$\\
29 & $119\;41\;\;4$ & $+127$ & $27\;32\;57$ & $+226$\\
30 & $119\;29\;26$ & $-50$ & $26\;58\;22$ & $+193$\\
31 & $119\;18\;36$ & $-306$ & $26\;23\;13$ & $+297$\\
April 1 & $119\;\;7\;38$ & $-382$ & $25\;50\;17$ & $+219$\\
\end{tabular}\end{center}

\SourceLeaf{347}\par\indent Für die erste Beobachtung von Herrn \textsc{Flaugergues} war in früheren Berichten eine um $10'$ größere Rektaszension angegeben. Dann wäre der Unterschied nur $259''$. Ich enthalte mich noch eines Urteils über die großen Abweichungen, bis ich einmal Zeit finde, selbst wieder Berechnungen zum Kometen anzustellen. Dazu komme ich derzeit noch nicht.


